\pdfoutput=1
\documentclass[11pt]{article}

\usepackage[utf8]{inputenc}
\usepackage[T1]{fontenc}
\usepackage{lmodern}
\usepackage[margin=1in]{geometry}
\usepackage{microtype}
\usepackage{booktabs}
\usepackage{tabularx}
\usepackage{array}
\usepackage{amsmath}
\usepackage{textcomp}
\usepackage[hidelinks]{hyperref}
\usepackage{url}
\usepackage[numbers,sort&compress]{natbib}

% Monospace inline code helper
\newcommand{\code}[1]{\texttt{#1}}

\title{One File, Many Readings: Cross-Extractor Content Divergence as
an Evasion Primitive in Email Attachment Pipelines}

\author{Ankitha Kundhuru}
\date{\today}

\begin{document}
\maketitle

\begin{abstract}
Production email-security pipelines extract the textual content of attachments
and score that text with downstream detectors (phishing/brand/malware
classifiers, increasingly LLM-based). These pipelines chain \emph{multiple}
content extractors across formats---Apache Tika, PDFBox, pypdf, pdfminer,
oletools. We show that the same attachment yields \emph{different extracted
content depending on which extractor runs}, and that this cross-extractor
divergence is an \emph{evasion primitive}: an attacker who knows (or probes) the
pipeline can craft an attachment whose extracted text scores benign while the
rendered document a human opens is a lure.

We contribute (1) an open-source \textbf{differential harness} measuring
divergence across $N$ extractors plus an OCR ``what the human sees'' oracle; (2)
a \textbf{vector taxonomy} beyond font remapping---invisible-layer,
text-as-image, image+invisible-decoy, optional-content, and
malformed-recoverable---implemented as programmatic mutations; and (3) a
\textbf{defense-gap} result: we re-implement the OCR font-verification defense
proposed by \emph{PDF Mirage} (USENIX~'17) and show it catches only the font
vector it was designed for, missing every non-font vector. On a real public
document corpus (GovDocs1), extractors diverge on 46\% of attachment-pair
comparisons, and two independent content detectors (a deterministic heuristic
and Claude) both flip malicious$\rightarrow$benign on 32\% of extractor pairs
for image-based masking vectors.
\end{abstract}

\section{Introduction}

An email-attachment detection pipeline does not score the \emph{bytes} of an
attachment; it scores the \emph{text an extractor pulls out of it}. That
extraction step is where security is quietly decided. Different extractors
implement different---individually spec-compliant---choices about encodings,
layer visibility, render modes, and recovery from malformed structure. The
result is that ``the content of this attachment'' is not a single string; it is
a function of which extractor you ask.

We frame this as an evasion primitive. If a pipeline's detector reads extractor
$X$'s output, an attacker crafts an attachment where $X$ reads benign text while
the rendered page---what the victim actually reads---carries the lure. This is
strictly more general than a single parser bug: it exploits \emph{legitimate
disagreement between compliant parsers}, so there is no single ``wrong'' parser
to patch.

\paragraph{Contributions.}
\begin{enumerate}
  \item \textbf{Measurement.} A differential harness (\S\ref{sec:harness}) that
  runs $N$ extractors plus an OCR render-oracle over an attachment and
  quantifies pairwise divergence, per-extractor blind-spot rates, and
  per-format/vector breakdowns. On real documents it measures a 46\% pairwise
  divergence rate (\S\ref{sec:measurement}).
  \item \textbf{Attack.} A vector taxonomy (\S\ref{sec:taxonomy}) and a
  detector-impact study (\S\ref{sec:detector}) showing which divergences flip a
  modern content detector's verdict---including an LLM detector (Claude).
  \item \textbf{Defense-gap.} A faithful re-implementation of PDF Mirage's
  font-verification defense and a demonstration that it does not generalize to
  the non-font vectors that dominate real attachments (\S\ref{sec:defense}).
\end{enumerate}

\section{Background and Related Work}

\paragraph{PDF Mirage: Content Masking Attack~\cite{pdfmirage}.} Weaponizes the
\emph{extracted $\neq$ rendered} gap via font remapping (a glyph renders as A
while \code{/ToUnicode} claims B) and proposes an OCR-based font-verification
defense. We differ on three axes: we are \emph{extractor-vs-extractor} (not
human-vs-single-parser), \emph{multi-format and multi-vector} (not font-only),
and we show their defense fails on our non-font vectors (\S\ref{sec:defense}).

\paragraph{Extract Me If You Can~\cite{extractme}.} Diffs a reference JavaScript
extractor (Adobe) against open-source extractors over 163k PDFs to evade malware
AV. Ours targets \emph{content} extraction (phishing/brand), uses a human-render
OCR oracle rather than an execute-oracle, and modern/LLM content detectors
rather than AV.

\paragraph{Body Obfuscation~\cite{bodyobfuscation}.} Measures email-\emph{body}
obfuscation prevalence (386 phishing emails, ten techniques) and its impact on
SpamAssassin/Rspamd scores. Ours is \emph{attachments}, a controlled automated
differential rather than a prevalence measurement, against modern detectors.

\paragraph{Positioning.} The intersection---a multi-extractor \emph{content}
differential in the email-attachment threat model, framed as a detector-evasion
primitive---is, to our knowledge, open.

\section{Threat Model}

The attacker can send an attachment and knows (or probes, via bounce/verdict
oracles) which extractor the target pipeline uses for a given format. The
defender runs one or more extractors and scores their output. The victim opens
the attachment in a normal renderer (the OCR oracle stands in for ``what the
victim reads''). The attacker's goal: a document scored benign by the pipeline's
extractor but read as a lure by the victim. We assume open-source tools and
public data throughout; no proprietary detection internals are used or required.

\section{The Differential Harness}
\label{sec:harness}

For one attachment the harness runs every applicable extractor and computes
pairwise similarity (normalized-text \code{SequenceMatcher} ratio;
\S\ref{sec:limitations} discusses stronger metrics), separating two axes:
\begin{itemize}
  \item \textbf{extractor-vs-extractor}---two parsers disagree on the same bytes.
  \item \textbf{extractor-vs-render}---a parser disagrees with the OCR oracle
  (render $\rightarrow$ pytesseract), i.e.\ reads something the human does not,
  or misses what they see.
\end{itemize}

Extractors: \code{pypdf}, \code{pdfminer.six}, Apache \code{Tika}, Apache
\code{PDFBox} (PDF and OOXML/OLE), \code{oletools} (OLE/macros), and
\code{ocr\_render} (PyMuPDF render + pytesseract) as the oracle. Every adapter
obeys a no-raise contract: a missing backend surfaces as a clean ``this
extractor was blind here,'' itself signal. A batch runner aggregates
divergence-rate statistics to machine-readable JSON/CSV.

\emph{Reproduce:} \code{python differ.py <file>};
\code{python matrix.py <corpus>}.

\paragraph{A note on Tika.} Tika bundles Tesseract and OCRs embedded images by
default, so it behaves as a hybrid parser+OCR and diverges from pure text
extractors on identical bytes (e.g.\ on a text-as-image PDF,
pypdf/pdfminer/pdfbox return empty while Tika returns ${\sim}0.66$ of the OCR
text). We treat Tika-OCR as a confound and recommend measuring both with and
without it (\code{OCR strategy NO\_OCR}).

\section{Divergence Vector Taxonomy}
\label{sec:taxonomy}

Each vector is a programmatic mutation that makes an attachment read one way to a
human and another to a text-layer parser (\code{mutations/}, validated
end-to-end by \code{mutations/validate.py}---all six produce their predicted
divergence).

\begin{table}[h]
\centering
\small
\begin{tabularx}{\linewidth}{@{}c l l X@{}}
\toprule
\# & Vector & Class & Mechanism \\
\midrule
1 & invisible-layer & text render mode 3 & decoy drawn invisibly; parsers read it, OCR/human don't \\
2 & text-as-image & rasterization & content is pixels, no text layer; parsers blind, OCR reads \\
3 & image+invisible decoy & 1+2 & lure as image, benign invisible decoy; parsers read benign \\
4 & optional-content & OCG layer OFF & decoy in a hidden PDF layer; renderers hide, parsers extract \\
5 & malformed-recoverable & corrupt xref & lenient parsers recover, strict ones error (availability divergence) \\
6 & font-encoding & glyph $\neq$ \code{/ToUnicode} & PDF Mirage primitive; parsers extract a scramble, OCR reads the glyphs \\
\bottomrule
\end{tabularx}
\caption{The six divergence vectors.}
\end{table}

Documented in \code{TAXONOMY.md}. Container/polyglot and email-multipart vectors
are enumerated as future extensions (the latter needs an \code{.eml}-level
harness).

\section{Measurement (D3)}
\label{sec:measurement}

\paragraph{Synthetic control.} On clean baseline PDFs, the pure-text extractors
agree exactly (pypdf vs pdfminer $=$ 0\% divergence). On the mutated corpus, 80\%
of files diverge---a clean separation confirming the vectors, not the harness,
drive divergence.

\paragraph{Real corpus.} We normalize public corpora with an \code{.eml}/\code{.mbox}
attachment loader. Apache SpamAssassin (2{,}400 spam emails $\rightarrow$ 31
attachments) is overwhelmingly image/HTML; document attachments are rare---itself
a finding about where document-parser evasion lives. For a document-rich sample
we use \textbf{GovDocs1} (Digital Corpora~\cite{govdocs1}), ${\sim}$1M
redistributable real \code{.gov} documents. On a 42-document sample (20 PDF / 10
DOC / 6 XLS / 6 PPT):

\begin{table}[h]
\centering
\small
\begin{tabular}{@{}l l@{}}
\toprule
metric & value \\
\midrule
overall pairwise divergence & \textbf{46.5\%} (66/142 comparisons) \\
files with $\geq$1 divergence & \textbf{69\%} \\
pdfminer vs pypdf & \textbf{55\%} (vs \textbf{0\%} on synthetic baseline) \\
pdfminer vs tika & 70\% $\cdot$ oletools vs tika 68\% $\cdot$ pdfbox vs pdfminer 50\% \\
\bottomrule
\end{tabular}
\caption{Real-corpus divergence (GovDocs1). \emph{Preliminary ($n=42$); the
harness scales to full GovDocs1 threads.}}
\end{table}

\paragraph{Worked example (\code{000816.pdf}, a US Census NAICS report).} The
same special-dash glyph is extracted four incompatible ways on identical bytes
(pypdf vs pdfminer similarity 0.25):

\begin{table}[h]
\centering
\small
\begin{tabular}{@{}l l r@{}}
\toprule
extractor & output for the glyph & chars \\
\midrule
pypdf & \code{/thrqtrEMdash} (internal glyph name leaked) & 31{,}614 \\
pdfminer & \code{(cid:1)} (unmapped CID) & 25{,}270 \\
tika & U+FFFD replacement char & 24{,}619 \\
pdfbox & \code{\textbackslash x01} (raw control byte) ---${\sim}$40\% of the text & 12{,}519 \\
\bottomrule
\end{tabular}
\caption{One glyph, four readings on identical bytes.}
\end{table}

This is a \emph{naturally occurring} font/encoding divergence in the wild---the
same class \S\ref{sec:taxonomy}'s vector 6 weaponizes deliberately.

\section{Detector Impact (D4)}
\label{sec:detector}

Divergence matters iff it changes a detector's decision. D4's scope is the five
non-font content-masking vectors (the font vector is evaluated separately as the
subject of the \S\ref{sec:defense} defense-gap experiment). For each such
masking vector we classify every extractor's output and the OCR ground truth
with two detectors---a deterministic phishing heuristic (offline, zero-cost,
reproducible) and an LLM detector (Claude, structured-output verdict). An
\textbf{evasion} is a (file, extractor) pair where the OCR truth scores malicious
but the extractor's text scores benign.

\begin{table}[h]
\centering
\small
\begin{tabular}{@{}l c c c@{}}
\toprule
detector & files w/ malicious truth & files evaded & extractor evasions \\
\midrule
heuristic & 5 & 2 & \textbf{6/19 (31.6\%)} \\
Claude (haiku) & 5 & 2 & \textbf{6/19 (31.6\%)} \\
\bottomrule
\end{tabular}
\caption{Detector-impact (evasion) results.}
\end{table}

The two independent detectors agree exactly---the evasion is a property of the
extraction divergence, not a heuristic artifact. Evasions concentrate on the
image-based vectors (text-as-image, image+decoy), where pypdf/pdfminer/pdfbox
read the benign decoy (or nothing) while the victim sees the lure. Notably Tika
does \textbf{not} evade there---its built-in OCR recovers the lure, closing the
blind spot the pure text extractors have (a concrete pipeline-configuration
consequence).

\section{Defense-Gap (D5)}
\label{sec:defense}

We re-implement PDF Mirage's OCR font-verification
(\code{defense/font\_verify.py}) as a faithful \emph{font-integrity} check: OCR
the rendered page and compare it to what the \textbf{visible} text layer claims
via \code{/ToUnicode} (visible $=$ \code{get\_texttrace} spans excluding
render-mode-3; OCG-hidden text is already absent). A page with no visible text
layer has no fonts to verify. Running it against every vector:

\begin{table}[h]
\centering
\small
\begin{tabular}{@{}l l c c l@{}}
\toprule
vector & class & font-based & flags? & verdict \\
\midrule
baseline & --- & no & no & true negative \\
\textbf{font\_remap} & font-encoding & \textbf{yes} & \textbf{yes} & \textbf{CAUGHT (as designed)} \\
invisible\_text & invisible-layer & no & no & MISSED \\
text\_as\_image & text-as-image & no & no & MISSED \\
image\_with\_decoy & text-as-image & no & no & MISSED \\
ocg\_hidden & optional-content & no & no & MISSED \\
malformed\_xref & malformed-recoverable & no & no & MISSED \\
\bottomrule
\end{tabular}
\caption{PDF Mirage's font-verification defense vs.\ every vector.}
\end{table}

\textbf{Font vectors caught 1/1; non-font vectors missed 5/5.} The defense is
single-extractor and font-specific: it verifies rendered glyphs against their
mapping. It cannot see content that is an image, an invisible/hidden layer, or a
structural-recovery divergence, and it says nothing about
extractor-vs-extractor disagreement---the primitive we measure. Since non-font
masking dominates real attachments, the accepted defense leaves the larger
surface open.

\section{Robustness Testing and Responsible Disclosure}
\label{sec:robustness}

Divergence is (mostly) spec-compliant behavior, not a memory-safety bug. To
probe for disclosable DoS/crash bugs we built a separate mutation-fuzzing
harness (\code{fuzz/}) that mutates real GovDocs1 seeds (radamsa) and runs each
parser in an isolated subprocess under a timeout + resource caps, classifying
crashes (signal death), hangs (timeout), and OOM/recursion as bugs distinct from
ordinary parse errors. A refined campaign (per target sequentially, 256\,KB
input cap, no CPU contention; 15k iterations on MuPDF, 6k each on pdfminer/pypdf,
4k on oletools, 1.5k each on Tika/PDFBox) rediscovered an
\textbf{unbounded-recursion condition in \code{olefile}'s directory handling}
(the OLE2 parser under \code{oletools}): a self-referential root directory entry
drives \code{OleDirectoryEntry.\_list()} into infinite recursion during
\code{listdir()} (minimized to a 1536-byte file). \textbf{This is not a novel
bug.} olefile issue \#103 (2018)~\cite{olefile103} reported infinite recursion
while \emph{opening} a malformed file; its fix moved the ``referenced more than
once'' guard to the top of \code{append\_kids()}, which stops the build-phase
recursion. Our case is a \emph{confirmed residual} of that fix: the root entry
is loaded directly and never passes through \code{append\_kids}, so it is never
marked ``used''; a self-referential root therefore slips into its own
\code{kids} and recurses in \code{\_list()} during \code{listdir()}---a path the
\#103 fix does not cover (cf.\ CWE-674~\cite{cwe674}). It is low-severity (a
catchable \code{RecursionError}; \#103 itself received no CVE), reproduces on the
current release, and is addressed by a one-line complement to the \#103 fix. We
report it, if at all, as an incomplete-fix hardening PR citing \#103, and make
\textbf{no CVE claim}. A weaker pdfminer.six super-linear-time observation
(73\,KB $\rightarrow$ ${\sim}$55 CPU-s) was also recorded. The takeaway is
methodological: the harness surfaced a real, minimizable robustness gap from real
seeds.

A methodological note worth reporting: a naive timeout-based hang detector
produces false positives from size-inflating mutators under CPU contention (an
earlier concurrent run flagged seven ``hangs'' that all completed on idle CPU).
Real DoS shows the opposite signature---\emph{small input, disproportionate
cost}---which the size cap plus per-run wall-time recording isolates. CVEs remain
upside, not a dependency of the contributions above.

\section{Limitations and Ethics}
\label{sec:limitations}

\begin{itemize}
  \item \textbf{Sample size.} Real-corpus numbers ($n=42$) and mutation-corpus
  numbers are preliminary; the harness scales to full corpora---this draft
  reports the phenomenon and the method, and larger runs are the natural next
  step.
  \item \textbf{Similarity metric.} \code{SequenceMatcher} is a first-pass
  metric; token-set or embedding distance may sharpen results without changing
  the qualitative story.
  \item \textbf{OCR oracle.} OCR is imperfect and slow on long documents (we
  expose a page cap and an off switch); it is a stand-in for human reading, not
  ground truth.
  \item \textbf{Ethics.} Open-source tools and public data only. Vectors are
  demonstrated on synthetic and public documents; any parser bug found via
  fuzzing is disclosed responsibly before publication. No proprietary detection
  internals are used.
\end{itemize}

\section{Conclusion}

The unit a content detector scores---extracted text---is not a property of the
attachment but of the extractor. We measured that cross-extractor divergence
directly (46\% on real documents), showed it flips modern content detectors
including an LLM (32\% of extractor pairs on image-masking vectors), and showed
the accepted content-masking defense is font-specific and misses the vectors
that dominate real attachments. The harness, taxonomy, and defense-gap
experiment are released for reproduction and extension.

\section*{Artifact Availability}

Open-source harness, mutation vectors, detectors, defense, and fuzzing harness:
\url{https://github.com/akundhuru/attachment-differ}, archived at Zenodo
(DOI: \href{https://doi.org/10.5281/zenodo.21894923}{10.5281/zenodo.21894923}).
Reproduction commands are inline per section; see \code{README.md},
\code{TAXONOMY.md},
\code{DEFENSE\_GAP.md}, \code{REAL\_CORPUS.md}, and \code{FUZZING.md}.

\bibliographystyle{plainnat}
\bibliography{refs}

\end{document}
